\documentclass[10pt,a4paper,twoside]{article}

% ============================================
% PAQUETES
% ============================================
\usepackage[utf8]{inputenc}
\usepackage[T1]{fontenc}
\usepackage[spanish]{babel}
\usepackage{geometry}
\usepackage{graphicx}
\usepackage{xcolor}
\usepackage{titlesec}
\usepackage{fancyhdr}
\usepackage{enumitem}
\usepackage{tabularx}
\usepackage{booktabs}
\usepackage{hyperref}
\usepackage{tcolorbox}
\usepackage{fontawesome5}
\usepackage{tikz}
\usepackage{caption}
\usepackage{tocloft}
\usepackage{setspace}
\usepackage{parskip}
\usepackage{lastpage}

% ============================================
% CONFIGURACIÓN DE PÁGINA (Compacta)
% ============================================
\geometry{
    top=1.5cm,
    bottom=1.5cm,
    left=1.5cm,
    right=1.5cm,
    headheight=15pt,
    columnsep=0.8cm
}

% ============================================
% COLORES KODEFY
% ============================================
\definecolor{kodefyBlue}{HTML}{2563EB}
\definecolor{kodefyDark}{HTML}{0F172A}
\definecolor{kodefyLight}{HTML}{E0E7FF}
\definecolor{kodefyGreen}{HTML}{10B981}
\definecolor{kodefyOrange}{HTML}{F59E0B}
\definecolor{kodefyRed}{HTML}{EF4444}
\definecolor{kodefyGray}{HTML}{64748B}
\definecolor{kodefyBg}{HTML}{F8FAFC}
\definecolor{tipBg}{HTML}{ECFDF5}
\definecolor{warnBg}{HTML}{FEF3C7}
\definecolor{noteBg}{HTML}{EFF6FF}
\definecolor{dangerBg}{HTML}{FEF2F2}

% ============================================
% ESTILOS DE TÍTULOS (Más compactos)
% ============================================
\titleformat{\section}
{\large\bfseries\color{kodefyDark}}
{\colorbox{kodefyBlue}{\textcolor{white}{\thesection}}}{0.5em}{}
[\vspace{-0.2em}\textcolor{kodefyBlue}{\rule{\columnwidth}{1.5pt}}]

\titleformat{\subsection}
{\normalsize\bfseries\color{kodefyDark}}
{\textcolor{kodefyBlue}{\thesubsection}}{0.5em}{}

\titleformat{\subsubsection}
{\normalsize\bfseries\color{kodefyGray}}
{\textcolor{kodefyBlue}{\thesubsubsection}}{0.4em}{}

\titlespacing*{\section}{0pt}{1.5ex plus 1ex minus .2ex}{1ex plus .2ex}
\titlespacing*{\subsection}{0pt}{1ex plus 1ex minus .2ex}{0.5ex plus .2ex}
\titlespacing*{\subsubsection}{0pt}{1ex plus 1ex minus .2ex}{0.5ex plus .2ex}

% ============================================
% CAJAS INFORMATIVAS (Compactas)
% ============================================
\newtcolorbox{tipbox}[1][Consejo]{
    colback=tipBg,
    colframe=kodefyGreen,
    coltitle=white,
    fonttitle=\bfseries\small,
    title={\faLightbulb\ #1},
    rounded corners,
    boxrule=1pt,
    left=4pt, right=4pt, top=2pt, bottom=2pt,
    before skip=6pt, after skip=6pt
}

\newtcolorbox{warningbox}[1][Importante]{
    colback=warnBg,
    colframe=kodefyOrange,
    coltitle=white,
    fonttitle=\bfseries\small,
    title={\faExclamationTriangle\ #1},
    rounded corners,
    boxrule=1pt,
    left=4pt, right=4pt, top=2pt, bottom=2pt,
    before skip=6pt, after skip=6pt
}

\newtcolorbox{notebox}[1][Nota]{
    colback=noteBg,
    colframe=kodefyBlue,
    coltitle=white,
    fonttitle=\bfseries\small,
    title={\faInfoCircle\ #1},
    rounded corners,
    boxrule=1pt,
    left=4pt, right=4pt, top=2pt, bottom=2pt,
    before skip=6pt, after skip=6pt
}

\newtcolorbox{dangerbox}[1][Cuidado]{
    colback=dangerBg,
    colframe=kodefyRed,
    coltitle=white,
    fonttitle=\bfseries\small,
    title={\faTimesCircle\ #1},
    rounded corners,
    boxrule=1pt,
    left=4pt, right=4pt, top=2pt, bottom=2pt,
    before skip=6pt, after skip=6pt
}

\newtcolorbox{stepbox}[1][]{
    colback=kodefyBg,
    colframe=kodefyBlue!30,
    coltitle=kodefyDark,
    fonttitle=\bfseries\small,
    title={#1},
    rounded corners,
    boxrule=0.5pt,
    left=4pt, right=4pt, top=2pt, bottom=2pt,
    before skip=4pt, after skip=4pt
}

% ============================================
% ENCABEZADOS Y PIES DE PÁGINA
% ============================================
\pagestyle{fancy}
\fancyhf{}
\fancyhead[LE,RO]{\textcolor{kodefyGray}{\small\thepage\ / \pageref{LastPage}}}
\fancyhead[LO]{\textcolor{kodefyBlue}{\small\bfseries KODEFY TECH}}
\fancyhead[RE]{\textcolor{kodefyGray}{\small\leftmark}}
\fancyfoot[C]{\textcolor{kodefyGray}{\footnotesize Manual de Usuario --- Versión 1.0}}
\renewcommand{\headrulewidth}{0.5pt}
\renewcommand{\headrule}{\hbox to\headwidth{\color{kodefyBlue}\leaders\hrule height \headrulewidth\hfill}}
\renewcommand{\footrulewidth}{0.3pt}
\renewcommand{\footrule}{\hbox to\headwidth{\color{kodefyLight}\leaders\hrule height \footrulewidth\hfill}}

% ============================================
% HYPERLINKS
% ============================================
\hypersetup{
    colorlinks=true,
    linkcolor=kodefyBlue,
    urlcolor=kodefyBlue,
    citecolor=kodefyBlue,
    bookmarksopen=true,
    pdfauthor={KODEFY TECH},
    pdftitle={Manual de Usuario - Sistema de Gestión KODEFY},
}

% ============================================
% TABLA DE CONTENIDOS
% ============================================
\renewcommand{\cftsecfont}{\bfseries\color{kodefyDark}\small}
\renewcommand{\cftsecpagefont}{\bfseries\color{kodefyBlue}\small}
\renewcommand{\cftsubsecfont}{\color{kodefyGray}\footnotesize}
\renewcommand{\cftsecdotsep}{\cftdotsep}
\setlength{\cftbeforesecskip}{2pt}

\begin{document}
\onecolumn

% ============================================
% PORTADA
% ============================================
\begin{titlepage}
\begin{tikzpicture}[remember picture, overlay]
    \fill[kodefyDark] (current page.south west) rectangle (current page.north east);
    \fill[kodefyBlue, opacity=0.3] (current page.north west) -- ++(0,-8) -- ++(\paperwidth,3) -- ++(0,5) -- cycle;
    \fill[kodefyBlue, opacity=0.15] (current page.north west) -- ++(0,-12) -- ++(\paperwidth,5) -- ++(0,7) -- cycle;
    \fill[kodefyBlue, opacity=0.2] (current page.south west) -- ++(0,4) -- ++(\paperwidth,-2) -- ++(0,-2) -- cycle;
    \fill[kodefyBlue, opacity=0.08] (14,-14) circle (6cm);
    \fill[kodefyBlue, opacity=0.05] (3,-20) circle (4cm);
\end{tikzpicture}

\vspace*{2cm}

\begin{center}
    {\fontsize{50}{60}\selectfont\textbf{\textcolor{white}{KODEFY}\textcolor{kodefyBlue}{TECH}}}
    
    \vspace{0.5cm}
    {\textcolor{kodefyBlue}{\rule{6cm}{2pt}}}
    
    \vspace{1.5cm}
    {\fontsize{24}{30}\selectfont\textcolor{white}{\textbf{Manual de Usuario}}}
    
    \vspace{0.3cm}
    {\Large\textcolor{kodefyBlue}{Sistema de Gestión Integral para Negocios}}
    
    \vspace{0.5cm}
    {\large\textcolor{gray!60}{Punto de Venta · Inventario · Reportes · Delivery}}
    
    \vspace{2cm}
    
    \begin{tcolorbox}[
        colback=kodefyBlue!10,
        colframe=kodefyBlue!40,
        width=10cm,
        arc=8pt,
        boxrule=1pt
    ]
    \centering
    \textcolor{white}{\large\bfseries Versión 1.0}\\[4pt]
    \textcolor{gray!70}{\small Octubre 2026 · Todos los derechos reservados}\\[4pt]
    \textcolor{gray!70}{\small Soporte: contacto@kodefytech.com}
    \end{tcolorbox}
\end{center}
\end{titlepage}

% ============================================
% TABLA DE CONTENIDOS
% ============================================
\newpage
\tableofcontents
\vspace{1cm}
\begin{center}
{\LARGE\bfseries\textcolor{kodefyDark}{¡Bienvenido a KODEFY TECH!}}
\end{center}
\vspace{0.5cm}
Estimado usuario, gracias por confiar en \textbf{KODEFY TECH} como la solución tecnológica para tu negocio. Este manual ha sido diseñado para guiarte paso a paso en el uso de todas las funcionalidades del sistema, desde la apertura diaria de tu caja hasta la generación de reportes detallados.

\begin{tipbox}[Sobre este manual]
Recomendamos leerlo completo la primera vez y después usarlo como referencia rápida cuando necesites recordar algún procedimiento.
\end{tipbox}

\newpage
\twocolumn
% ============================================
% INICIO DEL CONTENIDO EN DOS COLUMNAS
% ============================================

\section{Primeros Pasos}

\subsection{¿Qué es KODEFY TECH?}

\textbf{KODEFY TECH} es un sistema de gestión integral para negocios de comida. Funciona desde tu navegador web --- puedes usarlo desde una computadora, tablet o celular sin instalar nada.

Te permite manejar:
\begin{itemize}[itemsep=1pt, leftmargin=15pt]
    \item \textbf{POS:} Registrar pedidos rápidamente
    \item \textbf{Mesas:} Ver qué mesas están ocupadas
    \item \textbf{Cocina:} Pedidos en tiempo real
    \item \textbf{Inventario:} Control de stock
    \item \textbf{Delivery:} Gestión de repartos
    \item \textbf{Reportes:} Análisis detallado
    \item \textbf{WhatsApp IA:} Bot que atiende clientes
\end{itemize}

\subsection{¿Cómo accedo al sistema?}

Abre tu navegador web y ve a la dirección web proporcionada.

\subsubsection{Método 1: PIN de Negocio}

\begin{stepbox}[Paso 1: PIN]
Ingresa el \textbf{``PIN de Negocio (6 dígitos)''}. Cada negocio tiene su propio PIN único.
\end{stepbox}

\begin{stepbox}[Paso 2: Credenciales]
Una vez verificado el PIN, ingresa tu correo electrónico y tu contraseña personal.
\end{stepbox}

\subsubsection{Método 2: Multi-Negocio}

\begin{stepbox}[Varias sucursales]
\begin{enumerate}[itemsep=1pt, leftmargin=15pt]
    \item Presiona \textbf{``Tengo varios negocios''}.
    \item Ingresa tu correo y contraseña.
    \item Selecciona el negocio donde trabajarás hoy.
\end{enumerate}
\end{stepbox}

\subsection{Roles de Usuario}

\begin{itemize}[itemsep=2pt, leftmargin=15pt]
    \item \textcolor{kodefyBlue}{\textbf{Admin:}} Acceso total (crear usuarios, cambiar precios, reportes, configurar negocio).
    \item \textcolor{kodefyGreen}{\textbf{Cajero:}} Puede tomar pedidos, ver mesas, registrar ventas y ver el dashboard del día.
    \item \textcolor{kodefyOrange}{\textbf{Mozo:}} Puede tomar pedidos y asignarlos a mesas.
    \item \textcolor{kodefyRed}{\textbf{Repartidor:}} Ve únicamente sus pedidos de delivery con mapa GPS.
\end{itemize}

\begin{warningbox}
Solo el \textbf{Admin} puede crear usuarios y modificar la configuración.
\end{warningbox}

% ============================================
\section{Dashboard}

El \textbf{Dashboard} es tu panel de control principal.

\subsection{¿Qué muestra?}
\begin{itemize}[itemsep=1pt, leftmargin=15pt]
    \item \textbf{Ingresos:} Total vendido hoy.
    \item \textbf{Pedidos:} Cuántos atendiste hoy.
    \item \textbf{Stock:} Pollos y bebidas disponibles.
    \item \textbf{Gastos:} Compras registradas.
    \item \textbf{Top 5:} Productos más vendidos.
    \item \textbf{Métodos de pago:} Desglose de ingresos por efectivo, Yape, Plin, etc.
\end{itemize}

\begin{tipbox}
Se actualiza \textbf{en tiempo real}. Si otro cajero vende, los números se actualizan solos.
\end{tipbox}

% ============================================
\section{Apertura de Caja}

Es obligatorio abrir la caja antes de empezar a vender para indicar tu stock y dinero base.

\subsection{Proceso}
\begin{stepbox}[Cómo abrir]
\begin{enumerate}[itemsep=1pt, leftmargin=15pt]
    \item Ve a la sección \textbf{``Apertura''}.
    \item Ingresa los \textbf{Pollos enteros} disponibles.
    \item Ingresa \textbf{Kilos de papas} y \textbf{Litros de chicha}.
    \item Ingresa el \textbf{Dinero inicial} (base).
    \item Ingresa las \textbf{Bebidas} por marca y tamaño.
    \item Presiona \textbf{``Abrir Caja''}.
\end{enumerate}
\end{stepbox}

\begin{tipbox}[Carga automática]
Si cerraste caja correctamente el día anterior, el stock sobrante se \textbf{pre-carga automáticamente}.
\end{tipbox}

% ============================================
\section{Punto de Venta (POS)}

\subsection{Tipos de Pedido}
\begin{itemize}[itemsep=1pt, leftmargin=15pt]
    \item \textbf{Para Mesa:} Selecciona una mesa libre.
    \item \textbf{Para Llevar:} No se asigna mesa.
    \item \textbf{Delivery:} Pide dirección del cliente.
\end{itemize}

\subsection{Tomar un pedido}
\begin{stepbox}[Paso a paso]
\begin{enumerate}[itemsep=1pt, leftmargin=15pt]
    \item Elige el tipo de pedido.
    \item Agrega productos al carrito tocándolos.
    \item Si el producto tiene opciones (ej: parte del pollo), selecciónalas.
    \item Elige el método de pago (Efectivo, Yape, Plin, Tarjeta, Mixto).
    \item Presiona \textbf{``Confirmar Pedido''}.
\end{enumerate}
\end{stepbox}

\subsection{Impresión automática}
Al confirmar, se imprime el ticket para el cliente y la comanda para la cocina automáticamente (requiere configurar las IPs de las impresoras).

% ============================================
\section{Mesas y Cocina}

\subsection{Mesas}
\begin{itemize}[itemsep=1pt, leftmargin=15pt]
    \item \textbf{Verdes (Libres):} Listas para usar.
    \item \textbf{Rojas (Ocupadas):} Puedes tocarlas para agregar más productos, cambiar de mesa o cobrar.
\end{itemize}

\subsection{Pantalla de Cocina}
\begin{itemize}[itemsep=1pt, leftmargin=15pt]
    \item Los pedidos llegan con estado \textbf{``Pendiente''}.
    \item El cocinero ve detalles y notas.
    \item Al terminar, lo marca como \textbf{``Listo''} y el cajero/mozo recibe la notificación.
\end{itemize}

% ============================================
\section{Delivery y Radar}

\subsection{Pedidos a Domicilio}
Al crear un delivery, ingresas la dirección y el sistema calcula el \textbf{costo de envío} (si configuraste tu ubicación).

El estado avanza así: Buscando repartidor → Asignado → En camino → Entregado.

\subsection{Vista del Repartidor}
El repartidor tiene un \textbf{mapa GPS} en su celular para ver la ruta óptima de entrega y marcar el pedido como entregado.

% ============================================
\section{Historial y Anulaciones}

\subsection{Ventas del día}
Muestra todas las ventas. Puedes \textbf{Reimprimir tickets} o \textbf{Anular ventas}.

\begin{dangerbox}
\textbf{Anular} (solo Admins) cancela la venta, devuelve el stock y requiere ingresar un motivo para auditoría.
\end{dangerbox}

% ============================================
\section{Cierre de Caja}

Al finalizar el día, debes comparar lo registrado vs lo real.

\begin{stepbox}[Proceso de Cierre]
\begin{enumerate}[itemsep=1pt, leftmargin=15pt]
    \item Ve a \textbf{``Cierre de Caja''}.
    \item Revisa el resumen automático de ventas.
    \item Ingresa tu \textbf{conteo real}: pollos sobrantes, pollos usados por personal, mermas, bebidas sobrantes y dinero físico en caja.
    \item El sistema calculará diferencias (faltantes o sobrantes).
    \item Presiona \textbf{``Cerrar Caja''}.
\end{enumerate}
\end{stepbox}

\begin{warningbox}
Una vez cerrada la caja, no puedes registrar más ventas ese día.
\end{warningbox}

% ============================================
\section{Reportes y Gastos}

\subsection{Reportes}
Filtra por día o rango de fechas.
\begin{itemize}[itemsep=1pt, leftmargin=15pt]
    \item \textbf{Top Productos:} Gráfico de barras.
    \item \textbf{Métodos de pago:} Gráfico circular.
    \item \textbf{Ventas por hora:} Horarios pico.
\end{itemize}
Puedes exportar todos los datos a Excel.

\subsection{Gastos}
Registra tus compras diarias desde el Dashboard (botón ``Registrar Gasto''). Se restarán automáticamente del cierre de caja.

% ============================================
\section{WhatsApp con IA}

El bot atiende a tus clientes 24/7 por WhatsApp.

\textbf{¿Qué puede hacer?}
\begin{itemize}[itemsep=1pt, leftmargin=15pt]
    \item Responder dudas (horarios, menú).
    \item Tomar pedidos y calcular delivery.
    \item Gestionar reservas de mesas.
    \item Verificar vouchers de pago.
\end{itemize}

En la sección \textbf{``WhatsApp''} puedes apagarlo/encenderlo, cambiar su nombre, mensaje de bienvenida y configurar tus números de Yape/Plin.

% ============================================
\section{Inventario Avanzado}

Además de pollos y bebidas, puedes gestionar insumos completos (ej: arroz, aceite).
\begin{itemize}[itemsep=1pt, leftmargin=15pt]
    \item \textbf{Insumos:} Lista de materias primas.
    \item \textbf{Recetas:} Vincula insumos a productos para descontarlos automáticamente al vender.
    \item \textbf{Alertas:} Te avisa si hay stock bajo.
\end{itemize}

% ============================================
\section{Configuración}

Sección exclusiva para el \textbf{Admin}.

\begin{itemize}[itemsep=1pt, leftmargin=15pt]
    \item \textbf{Productos:} Cambia precios y actívalos.
    \item \textbf{Usuarios:} Crea empleados y asigna roles.
    \item \textbf{Mesas:} Crea mesas por piso.
    \item \textbf{Negocio:} Cambia nombre, logo y colores.
    \item \textbf{Impresoras:} Pon las IPs de las térmicas.
    \item \textbf{Catálogo de Bebidas:} Agrega o quita marcas.
\end{itemize}

% ============================================
\section{Tips Extra}

\subsection{Instalar como App}
KODEFY TECH se puede instalar: abre Chrome en tu celular, toca los tres puntos y selecciona \textbf{``Agregar a pantalla de inicio''}. Funcionará como una app nativa a pantalla completa.

\subsection{Tema Oscuro}
Usa el botón de la luna arriba a la derecha para activar el modo oscuro (ideal para la noche o ahorrar batería).

\vspace{1cm}

\begin{center}
    {\large \faWhatsapp\ Soporte: contacto@kodefytech.com}
\end{center}

\end{document}
